\section{Corrections to versions 1--9}\label{sec:corrections}

This section lists the principal claims that are withdrawn or changed, so that a reader of an
earlier version can see which statements should no longer be relied on.  It is not guaranteed
complete: where an earlier version and this one disagree, this one is intended.  Items marked \emph{wrong} are
contradicted by a computation or an argument in this version; items marked \emph{not
re-verified} may be right but were not recomputed with the exact methods of
Section~\ref{sec:methods} and are therefore no longer claimed.

\paragraph{Withdrawn as wrong.}
\begin{enumerate}[label=C\arabic*.,leftmargin=2.6em]
\item \emph{Golden-ratio minimum 52.}  A 51-ray uncolorable subset exists
(Section~\ref{sec:catalogue}).
\item \emph{``OCUS-certified'' minima.}  The certification code was unsound
(Section~\ref{sec:methods}).  All minima in this version were recomputed with a different
procedure.  The values 31, 33, 33, 33 and 43 for the integer, Peres, $\sqrt{-2}$, Eisenstein and
$\Q(\sqrt{-7})$ raw pools survive; 52 does not.
\item \emph{``Islands are isolated; no continuous path connects one KS-supporting alphabet to
another.''}  The alphabet $\{0,\pm1,\pm x,\pm x^{*}\}$ is uncolorable along the whole circle
$|x|^2=2$, which joins the Peres, $\sqrt{-2}$ and $\Q(\sqrt{-7})$ alphabets
(Theorem~\ref{thm:circle}).  The sweep of 164 angles offered as evidence could not have
established isolation in any case.
\item \emph{The norm bound}: ``generators with $|x|^2\ge3$ that are not roots of unity produce no
KS sets.''  It fails for $x=2$, the alphabet of $\CK$; $|x|^2$ is not an invariant of the pool
($x$ and $1/x$ give the same rays); and version~9 itself reported uncolorable pools from
generators of squared modulus 4, 5 and~6.
\item \emph{The two- and three-mechanism classifications} (``modulus-2'', ``phase'',
``minimal-polynomial'', ``indirect modulus-2'').  The statement that an uncolorable pool requires
a cancellation identity is true by definition.  The list of mechanisms was extended after each
counterexample and made no prediction.  What survives is Theorem~\ref{thm:twosymbol},
Theorem~\ref{thm:6n}, Theorem~\ref{thm:circle} and Table~\ref{tab:zerosums}.
\item \emph{Completeness of the table of three-term zero-sums.}  The earlier table was compiled
by inspection, its justification miscounted the cases, it omitted two-term sums, and it described
solution sets that are circles or lines as isolated points.  Table~\ref{tab:zerosums} is the
complete enumeration.  The accompanying claim that an SMT solver had certified completeness
overstated what was run.
\item \emph{``Universal merge saturation''} and the conjecture built on it.  Identifying two
non-orthogonal rays of any uncolorable configuration gives an uncolorable configuration: a
coloring of the quotient pulls back to the original by giving the two rays the same value.  The
3{,}756 SAT calls verified this lemma; nothing about minimal KS sets follows.  (The printed
definition also described deletion of a ray, not identification.)
\item \emph{The Galois-theoretic reformulation.}  The statement that $(\Z/n\Z)^{*}$ surjects onto
$\Z/6\Z$ if and only if $6\mid n$ is false in both directions ($n=6$, $n=7$); the conductor of
$\Q(\zeta_6)$ is 3, not 6; the norm characterization fails for real quadratic generators; and the
list of ``KS-supporting'' imaginary quadratic fields contradicted the paper's own tables.  The
parallel drawn with Gleason's theorem and with the matrix trace is withdrawn with it.
\item \emph{The Hoffman-bound explanation of CSW advantages}, the ``auxiliary ray'' mechanism and
the description of the Peres pool as disconnected (Section~\ref{sec:csw}).
\item \emph{``$|S_A|+|S_B|\ge|\mathcal B|$ in every case''} (Section~\ref{sec:cabello}), and the
claim that every Bell scenario obtainable from $\CK$ is obtainable more cheaply from the
Eisenstein set, which rested on 200 heuristic trials.
\item \emph{``The Peres and $\Z[\sqrt{-2}]$ islands are isomorphic.''}  True of the 33-ray minimum
sets, false of the pools (Computed fact~\ref{fact:pools}).
\item \emph{Counts.}  $\{0,\pm1,\pm i\}$ has 31 rays, 69 pairs, 7 triads, not 25/48/6.  The
completed pools of $\sqrt3$ and $\sqrt5$ have 133 rays, not 145.  The Eisenstein pool is not
closed under completion (165 rays after completion, not 57).  The $\Q(\sqrt{-7})$ pool has 522
orthogonal pairs in total; version~9 said its 76 auxiliary rays ``add'' 522.  The integer pool
has 6 minimum sets in one symmetry class, not 162 in six.
\item \emph{The Hermitian cross-product formula} was printed with a spurious alternating sign;
the third ray of an orthogonal pair $v,w$ is $[(v\times w)^{*}]$.
\item \emph{``All coordinates are exact, with no floating-point approximation''} was not true of
all the scripts cited.  In this version it is true of the pool counts, colorability decisions,
minima, isomorphism tests and bipartite costs; the semidefinite programs, attained values and
rigidity ranks are numerical.
\item \emph{Further items.}  The claim that enlarged Eisenstein alphabets keep minimum 33 (once
the letters $\pm2$ are admitted the pool contains $\CK$).  The use of ``contextual fraction'' for
a ratio of bases, a term with an established different meaning.  The old certification script
could print ``certified'' after finding a smaller counterexample.  The statement that the
bipartite search was exhaustive.
\item \emph{Smaller points.}  The 13-ray pool was called too small ``because the minimum known is
31'': 31 is an upper bound, and the relevant fact is the lower bound 24.  Infinitesimal rigidity
was called necessary for global rigidity; it is neither necessary nor sufficient.  The rigidity
theorem of Trandafir and Cabello concerns realizations of the $\CK$ graph and does not imply that
every 31-ray KS set is $\CK$.  The Kirchweger--Peitl--Szeider reference had the wrong venue,
pages and an author's initial.
\end{enumerate}

\paragraph{No longer claimed (not re-verified).}
The survey of imaginary quadratic fields of class number 2--4; the survey of 28 two-generator
alphabets; the quartic alphabets; mixed alphabets with two independent irrationals; the
trigonometric sweep; the icosahedral directions; the quaternionic and octonionic alphabets (for
the latter the notions of ray and orthogonality were never defined); the comparison with the
97-ray closure of Trandafir and Cabello; the $N(S)$ comparison with Cortez, Morales and Reyes;
the ``graph universality'' of $\CK$ across group-orbit constructions; the ``30{,}000 random
hypergraphs''; the identification of the coordinate field of the original 117-ray set of Kochen
and Specker; the statement that the 43-ray and golden-ratio sets are new to the literature.

\section{Open questions}\label{sec:limitations}

\begin{enumerate}
\item \emph{Other loci.}  Table~\ref{tab:zerosums} has two circles and three lines.  For the
two-symbol alphabet all of them are settled (Theorem~\ref{thm:twosymbol}): none is uncolorable
away from isolated points.  For the conjugate-closed alphabet we have settled the circle
$|x|^2=2$ only.  What happens along $\mathrm{Re}\,x=-\tfrac12$ (which contains $\omega$) and along
the unit circle when $x^{*}$ is included is not known.  (On the imaginary axis $x^{*}=-x$, so
the conjugate-closed alphabet is the two-symbol one and that case is settled.)
\item \emph{Three symbols.}  The method of Theorem~\ref{thm:twosymbol} applies in principle to
$\{0,\pm1,\pm x,\pm y\}$: the relations are again finitely many, but the parameter space has
four real dimensions and we have not attempted the enumeration.  Could a three-symbol pool
contain a KS set of fewer than 31 rays?
\item \emph{The second flex.}  Two infinitesimal flexes of the 45 words were found numerically
at both phases tested (Table~\ref{tab:rigidity}).  Is the second one a genuine deformation, leading off the
family of Theorem~\ref{thm:circle}?
\item \emph{All minimum sets on the circle.}  We exhibit one 45-word minimum set; the others have
not been enumerated, nor has the set of phases at which the pool gains orthogonalities.
\item \emph{Completion.}  Is there a finiteness criterion for completed pools?  Why does
completion make the golden-ratio and $\sqrt[3]2$ pools uncolorable but not those of $\sqrt3$ and
$\sqrt5$?
\item \emph{Uncertified minima.}  The entries of Table~\ref{tab:catalogue} marked ``$\le$''.
\item \emph{The gap.}  This paper does not close the gap between the lower bound 24 and
the 31 rays of Conway and Kochen.  Theorem~\ref{thm:twosymbol} excludes a smaller set with a
two-symbol alphabet and nothing more; the other minima are certified within named finite pools
only.  Earlier
versions suggested that the alphabet analysis could prune the hypergraph searches
of~\cite{LiBrightGanesh2024,KirchwegerPeitlSzeider2023}; we see no basis for that, since a
hypothetical smaller set need not have a small alphabet.
\item \emph{Catalogue comparison.}  Whether the 43-, 45- and 47-ray sets appear in the automated
catalogues of three-dimensional KS sets.
\end{enumerate}

\section{Conclusion}\label{sec:conclusion}

Reading the coordinates off Peres's two cubes suggests a one-parameter question: for which $x$
does the alphabet $\{0,\pm1,\pm x\}$ generate an uncolorable pool?  The answer is now complete
(Theorem~\ref{thm:twosymbol}).  Up to the sign and inversion of $x$, which do not change the pool, it is the two classical cases
$x=\sqrt2$ and $x=2$ and nothing else for real $x$; for complex $x$ it adds only $\sqrt{-2}$ and
the primitive third and sixth roots of unity; and in no case is there a KS set of fewer than 31 rays.  One unit and one
further number along each axis cannot improve on Conway and Kochen.  The more interesting
structure appears when the alphabet is enlarged.  With the conjugate
included, the pool is uncolorable on the entire circle $|x|^2=2$ and colorable off it for
generic~$x$; twelve orthogonalities of the form $1+1-|x|^2=0$ make the difference; the Peres set
and its imaginary twin are the two degenerate points of a 45-ray family; and the 43-ray set over
$\Q(\sqrt{-7})$ sits at a point of the circle where one more relation holds.  For the
roots-of-unity alphabets the answer is exact as well: uncolorable if and only if $6\mid n$, with
the primes 2 and 3 doing different jobs.

The earlier versions of this paper claimed more than this and were wrong in ways that matter.
We have tried to make this version checkable instead: every number is produced by a named script
in exact arithmetic or is labelled numerical, every minimum is either certified by a procedure
whose soundness is argued in two paragraphs or reported as a bound, and the withdrawn claims are
listed.

\subsection*{Reproducibility}

Code and data: \url{https://github.com/michaelkernaghan/contextuality}, release tag
\texttt{v10.1}.  Scripts for this
version, with the result files they write under \texttt{results/}:

\begin{itemize}\setlength{\itemsep}{0pt}
\item \texttt{ks\_recertify\_v10.py}: pools of Table~\ref{tab:catalogue}, with certified minima
and witnesses, written to \texttt{recertify\_v10.jsonl}.  It uses the exact fields, pools and
completion of \texttt{ks\_verify\_large\_minima.py} and the certifier
\texttt{ks\_certify\_fast3.py}.
\item \texttt{ks\_tables\_v10.py}: Tables~\ref{tab:alphabets} and~\ref{tab:heegner} and the
four-symbol and completed pools quoted in the text (\texttt{tables\_v10.jsonl}).
\item \texttt{ks\_zero\_sums\_v10.py}: Table~\ref{tab:zerosums}.
\item \texttt{ks\_two\_symbol\_theorem.py}: Theorem~\ref{thm:twosymbol} and
Table~\ref{tab:twosymbol} (\texttt{two\_symbol\_theorem.json}, with a witness for every
uncolorable case); about five seconds.
\item \texttt{ks\_verify\_cyclotomic\_theorem.py}: Table~\ref{tab:cyc} and the structural
statements of Theorem~\ref{thm:6n}, $n\le30$.
\item \texttt{ks\_modulus2\_circle.py}: Theorem~\ref{thm:circle}, Table~\ref{tab:words},
Appendix~\ref{app:triads}.
\item \texttt{ks\_extras\_v10.py}: Table~\ref{tab:profiles}, the isomorphism tests and the
enumeration of all minimum sets (\texttt{extras\_v10.jsonl}).
\item \texttt{ks\_bpqs\_exact.py}: Table~\ref{tab:bpqs}.  \texttt{ks\_csw\_exact.py}:
Table~\ref{tab:csw}.  \texttt{ks\_rigidity\_v10.py}: Table~\ref{tab:rigidity}.
\item \texttt{make\_alphabet\_cube\_figures.py}: Figs.~\ref{fig:triptych} and~\ref{fig:generic}.
\end{itemize}

Software: Python~3.13, PySAT (Glucose4, RC2), NetworkX (VF2), NumPy, SciPy (HiGHS), CVXPY
(Clarabel).  Scripts from earlier versions remain in the repository and should not be used to
reproduce this one.

\subsection*{How this version was produced}

Computations, proofs and text were produced with a large language model (Claude, Anthropic)
under the author's direction.  Version~9 was then reviewed in full by that model and, separately,
by a second model (GPT, OpenAI) instructed to act as a hostile examiner with access to the
repository; the two reviews agreed on the main defects and each found errors the other missed.
The present version was written after those reviews.  Agreement between language models is not a
proof.  The safeguards we rely on are the ones listed in the Conclusion: exact arithmetic, a
certifier with a short soundness argument, positive controls (the bipartite search reproduces the
published $5\times9$ and $7\times9$; the pool counts reproduce those of earlier versions where
those were right), and an explicit list of what was not checked.  Theorem~\ref{thm:6n} has a
proof that can be read in ten minutes, and independent checking of it and of
Theorems~\ref{thm:twosymbol} and~\ref{thm:circle} is invited.

\subsection*{Acknowledgments}

The author thanks Ad\'an Cabello for pointing out the connection between the Eisenstein alphabet
and his 33-ray set; Manuel Flores Gordillo for finding the error in the rigidity computation of
versions 1--7; Jorge Lucas Gonzalez for the supergolden construction and an exact verifier;
Markus Frembs for the trace-zero relation and for engaging with the earlier classification; and
Curtis Bright for a bibliographic correction.  The author accepts responsibility for the content,
including the errors of the earlier versions.
